\documentclass[11pt]{article}
\usepackage{iftex}
\usepackage[T1]{fontenc}
\usepackage{lmodern,microtype}
\usepackage[a4paper,margin=25mm,top=23mm,bottom=24mm]{geometry}
\usepackage{amsmath,amssymb,amsthm,mathtools}
\usepackage{booktabs,array,enumitem}
\usepackage{xcolor}
\usepackage[colorlinks=true,linkcolor=blue!50!black,citecolor=blue!50!black,urlcolor=blue!50!black]{hyperref}

\ifdefined\pdftrailerid\pdftrailerid{}\fi
\setlength{\emergencystretch}{1.5em}
\numberwithin{equation}{section}

\newtheorem{theorem}{Theorem}[section]
\newtheorem{proposition}[theorem]{Proposition}
\newtheorem{lemma}[theorem]{Lemma}
\newtheorem{corollary}[theorem]{Corollary}
\theoremstyle{definition}
\newtheorem{definition}[theorem]{Definition}
\theoremstyle{remark}
\newtheorem{remark}[theorem]{Remark}

\newcommand{\R}{\mathbb{R}}
\newcommand{\Q}{\mathbb{Q}}
\newcommand{\Z}{\mathbb{Z}}
\newcommand{\N}{\mathbb{N}}
\newcommand{\dd}{\,\mathrm{d}}
\newcommand{\pos}[1]{[#1]_+}
\newcommand{\Bta}{\mathrm{B}}
\DeclareMathOperator{\supp}{supp}
\DeclareMathOperator{\sgn}{sgn}
\DeclareMathOperator{\ord}{ord}
\newcommand{\Hhyp}{H_{\mathrm{hyp}}}
\newcommand{\PP}{\mathcal{P}}
% Print the companion's number on paper. The PDF checker matches #1 to its LaTeX label.
\newcommand{\othernumber}[2]{#2}
% Sibling PDF, PDF destination, source label, printed number.
\newcommand{\paperlink}[4]{\href{#1.pdf\##2}{#4}}


\ifdefined\pdfinfoomitdate\pdfinfoomitdate=1\fi
\ifdefined\pdfsuppressptexinfo\pdfsuppressptexinfo=15\fi
\newcommand{\disk}{\mathbb D}
\newcommand{\ev}{\operatorname{ev}}
\newcommand{\cL}{\mathcal L}
\newcommand{\cN}{\mathcal N}
\newcommand{\F}{\mathcal F}
\setlist{topsep=4pt}
\makeatletter\let\inputrows\@@input\makeatother

\newcommand{\splitcount}[2]{\left(#1,#2\right)_{\mathrm{terms}}}
\newcommand{\Mclr}{M_{\mathrm{clr}}}

\hypersetup{pdftitle={Degree-capped fields and finite return certificates},pdfauthor={Haibo Lu},filecolor=blue!50!black}
\renewcommand{\thesection}{S\arabic{section}}
\renewcommand{\thetable}{S\arabic{table}}
\setcounter{tocdepth}{1}
\makeatletter
\let\basesectiontoc\l@section
\renewcommand{\l@section}[2]{\begingroup\let\orig@numberline\numberline\renewcommand{\numberline}[1]{\makebox[2.5em][l]{##1}}\basesectiontoc{#1}{#2}\endgroup}
\makeatother
\title{Degree-capped fields and finite return certificates\\[4pt]\large Supplement to \emph{A superquadratic lower bound for limit cycles of planar polynomial vector fields with \(N\) monomials and degree at most \(N-1\)}}
\author{Haibo Lu}
\date{November 2026}
\begin{document}
\maketitle
\begin{abstract}
We give the literature and ordinary-support calculations behind the degree-capped comparisons. A one-coordinate lift gives four hyperbolic cycles at five terms and degree four. Two-well fields cover six through eight terms, and a fully specified rational field has 783 hyperbolic cycles at degree 31. We also evaluate the fold loss at a dyadic budget and exhibit both an attainable and an obstructed primitive on a fifth-degree fold.
\end{abstract}
The source archive contains the integer and interval calculations \cite{LuReplay}. Separate companion notes treat fields of unrestricted degree and returns on a single-row binary saddle tree.

\tableofcontents
\section{Literature comparisons at exact support}\label{sup:sources}
Numbered results, equations and tables in this supplement carry an S prefix. References to the main article retain its numbering. Write \(N=|\supp P|+|\supp Q|\) after ordinary expansion, combination of equal powers in each component, and removal of zero coefficients. In \(\splitcount{a}{b}\), \(a\) and \(b\) are the term counts in \(P\) and \(Q\), respectively. We use \(D=\max(i+j)\) among the surviving pairs in either component. Gasull and Santana's Theorem~2 and Table~1 give \(12,12,12,16,20,24,32\) isolated cycles at \(N=4,\ldots,10\), with unrestricted degree \cite[Theorem~2 and Table~1]{GasullSantana2024}. Their Theorem~1 gives \(N^2/2-3N-8\) for every \(N\ge9\). At \(N\le3\) the sharp zero bound follows from Buzzi, Carvalho and Gasull \cite{BuzziCarvalhoGasull2021}. For four terms their proof uses the smooth family \((y,-x+a_0y+a_1y^3)\). First-order averaging supplies two nonzero coefficients giving a hyperbolic cycle. This accounts for the four-term entry at \(D=3\).

\paragraph{A weak focus before and after squaring.}
With \(k=1,\ldots,5\), write Gasull and Santana's (21) \cite[pp.~11--12]{GasullSantana2024} as
\[
 \dot x=a_1xy^2+\sum_{j=2}^k a_jy^{j+1}-a,\qquad
 \dot y=xy-1,\qquad a=\sum_{j=1}^k a_j .
\]
For \(k=4,5\) the construction unfolds a weak focus at \((1,1)\). For \(k=5\), the author version of \cite[(21), p.~11]{GasullSantana2024} (arXiv:2405.04281v3) prints \(a_2=-161/17\). In the coordinates of Section~\ref{sup:focus} this gives \(c_2=19803/7162\). Holding the other displayed coefficients fixed, \(c_2=0\) gives \(a_2=-161/99\). Here \(c_2=c_4=c_6=c_8=0\), \(c_{10}=-1/27\), and the Jacobian of \((c_2,c_4,c_6,c_8)\) with respect to \((a_2,a_3,a_4,a_5)\) has determinant \(33/2\). Together with \(c_0=1+a_1\), these give the order and independent unfolding directions used in the five-cycle argument.

The substitution \(x=u^2,\ y=v^2\) and the time factor \(2uv\) give (22), with components \(vP(u^2,v^2)\) and \(uQ(u^2,v^2)\). On each open quadrant the substitution is invertible and the time factor is nonzero. The last paragraph of their proof of Theorem~2 \cite[Theorem~2, Eq.~(22)]{GasullSantana2024} states that for \(k=1,2,3\) there are parameters yielding \(k+3\) terms and at least \(k\) cycles per quadrant. Their four-term field (3), p.~4, already gives four cycles by Hopf bifurcation. For \(k=2\), (22) has component counts \(\splitcount{3}{2}\), degree seven and eight cycles. For \(k=3\) the truncation sets \(a_4=a_5=0\) and gives
\[
 \dot u=a_3v^9+a_2v^7+a_1u^2v^5-av,\qquad \dot v=u^3v^2-u .
\]
Here \(a=a_1+a_2+a_3\) and all four coefficients \(a,a_1,a_2,a_3\) are nonzero. Thus the component counts are \(\splitcount{4}{2}\), the degree is nine, and at least twelve cycles coexist for a suitable parameter choice. Gasull and Santana give an existential parameter choice for \(k=3\). Simple-return realizations for all three small truncations are derived in Section~\ref{sup:focus}.

For \(N=4,5,6\) Section~\ref{sup:focus} proves that the following lower counts are hyperbolic at suitable parameters. Each entry is \((C,D)\) after ordinary expansion, with component counts \(\splitcount{k+1}{2}\).
\begin{center}\small
\begin{tabular}{rccc}\toprule
\(N\)&Original (21)&One square&Two squares (22)\\\midrule
4&\((1,3)\)&\((2,4)\)&\((4,7)\)\\
5&\((2,3)\)&\((4,4)\)&\((8,7)\)\\
6&\((3,4)\)&\((6,4)\)&\((12,9)\)\\\bottomrule
\end{tabular}
\end{center}
For \(k=4,5\) the original isolated-cycle counts are \(4,5\) at degrees \(5,6\); after the two squares they become \(16,20\) at degrees \(11,13\).

\paragraph{Cubic comparisons and the Songling quadratic.}
Gasull and Santana attribute the five-term count of three to Chow, Li and Wang, \S4.2 \cite{GasullSantana2024,ChowLiWang1994}. The degree of the five-term field in Table~\othernumber{tab:small}{3} of the main article comes from Chen and Chen's independent equation
\[
 \dot x=y-a_1x-a_2x^3,\qquad \dot y=a_3x+a_4x^3 .
\]
In the interior of region~IX there are two small unstable cycles and a large stable one. With all four coefficients nonzero the component counts are \(\splitcount{3}{2}\) and \(D=3\) \cite[Table~2]{ChenChen2018}. Section~\ref{sup:quartic} gives a further cubic realization at \(n=1,d=2\).

The cuspidal-loop and figure-eight-loop studies of Dumortier and Li give cubic systems with at least four and five cycles, respectively \cite{DumortierLi2,DumortierLi4}. In their family \(\dot x=y,\ \dot y=P_3(x)+yQ_2(x)\), the subscripts specify polynomial degrees, and the Hamiltonian of degree four gives a vector field of degree three. The entries marked \(^{a}\) in Table~\othernumber{tab:small}{3} of the main article, at exact term counts \(N=6,7\), follow Gasull and Santana's attribution to these papers in their introduction \cite{GasullSantana2024}. The entries record isolated cycles at the family level.

Li, Liu and Yang's nine-term cubic has \(H_x=x(x+1)(x-\lambda)\), \(H_y=y^3-k^2y\), with a perturbation \(y(\alpha_1+\alpha_2 x+\alpha_3 x^2+\alpha_4 y^2)\) in the second component. The signs and simple zeros in Proposition~1 give thirteen isolated cycles \cite[pp.~3610--3613, Eqs.~(3)--(4)]{LiLiuYang}. Opposite signs on a thirteenth bracket allow a nearby choice occupying all nine positions, even if a coefficient initially vanishes.

For the eight-term comparison, Galias and Tucker use the Songling field
\[
 \dot x=\lambda x-y-10x^2+(5+\delta)xy+y^2,\qquad
 \dot y=x+x^2+(-25+8\varepsilon-9\delta)xy
\]
with \(\delta=-10^{-13},\ \varepsilon=-10^{-52},\ \lambda=-10^{-200}\). All five coefficients in the first component and all three in the second are nonzero. Theorem~1 gives exactly four cycles for this fixed quadratic. Lemma~3 encloses the four return derivatives in intervals strictly to either side of one, establishing hyperbolicity of each cycle \cite[Eq.~(1) and Lemma~3, pp.~2, 4]{GaliasTucker2022}.

\paragraph{The fields in Gasull and Santana's Proposition 1.}
Gasull and Santana use the simple roots from their Lemma~1 in Proposition~1 \cite{GasullSantana2024}. Their derivative identity (7) corresponds to the polynomial
\[
 Q_r(x)=\prod_{k=-r}^r(x-k)=x\prod_{j=1}^r(x^2-j^2).
\]
In the accepted manuscript and arXiv version~3, (6) prints an additional \(x\) to the left of the first product. At the origin the literal expression has a multiple zero and fails (7). The displayed simple-root choice has \(r+1\) nonzero odd powers by elementary symmetric sums, \(Q'_r(j)=(-1)^{r-j}\prod_{k\ne j}|j-k|\), and alternating centers on \(y=\pm1\) for \(\dot x=y-y^3,\dot y=Q_r(x)\). These are the centers used in the proof of the proposition. With \(m_0=0,\ c_0=1\) use the perturbation in that proof,
\[
 P=y-y^3+\epsilon\sum_{i=0}^n c_i x^{2(n-i)+1}y^{2m_i},\qquad Q=Q_r(x).
\]
Successively large integers \(m_i\) localize the \(i\)-th contribution between the chosen ovals, preserving finitely many strict previous signs. The coefficients \(c_i\) are nonzero with alternating signs and include the normalization of their vertical marks. Pure \(y\) terms cannot collide with the \(n+1\) terms having distinct positive \(x\) exponents. Every factor of \(Q_r\) is \(x-j\), and the odd part in \(y\) of \(P(j,y)\) is \(y-y^3\), so the two components are coprime. Thus \(N=n+r+4\) and
\[
 G(n,r,\mathbf m)=\max\!\left(3,\,2r+1,\,\max_{0\le i\le n}\{2(n-i)+1+2m_i\}\right).
\]
The construction guarantees at least \(2n(r+1)+n(1+(-1)^r)\) odd crossings. To maximize this count at \(N=n+r+4\), write \(K=N-2\). At \(r=2t-1\) it is \(4t(K-2t-1)\), below \(4t(K-2t)\) at \(r=2t-2\). The latter equals \(K^2/2-8(t-K/4)^2\). Rounding \(t\) to the nearest integer gives \(4\lfloor K^2/8\rfloor\) for \(N\ge9\), an optimization within the proposition. In the proof of Theorem~2 the choices \((n,r)=(3,2),(4,2)\) give 24 and 32, respectively \cite[Proposition~1 and proof of Theorem~2]{GasullSantana2024}. Since \(0=m_0<m_1<\cdots<m_n\) are integers, \(m_i-i\) is nondecreasing. The largest perturbation degree is therefore \(2m_n+1\). For these two choices it also exceeds \(2r+1=5\), so write \(G_3=2m_3+1\ge7\) and \(G_4=2m_4+1\ge9\) for their actual degrees. For these choices \(n\in\{3,4\}\) and \(N=n+6\), so the cap forces \(m_n\le\lfloor(n+4)/2\rfloor=n\). The identity with \(n\) holds precisely for \(n=3,4\). Strict increase then forces \(m_i=i\) throughout. The successive large-exponent choices prove the sign brackets without furnishing either of these specializations. The nine-term quartic discussed below is a separate construction.

At \(N=9\) the stationary count of the main paper is \(25\), compared with \(24\) from the optimized formula. For \(N\ge10\), twice their difference is at least \((N-6)^2-10>0\). The four-quadrant examples for \(N=4,5,6\) in Gasull and Santana start from the three-cycle planar S-system of Boros and Hofbauer \cite[\S7]{BorosHofbauer2019}.

There is a same-support hyperbolic refinement. On any one regular annulus choose the orientation of its area integral and vary the existing coefficient of \(x^{2n+1}\) by \(c\) in the perturbation. The first-order change is \(cJ\), with \(J(h)=\iint_{\Omega_h}(2n+1)x^{2n}\dd x\dd y>0\), up to the same orientation. The endpoints of each bracket retain their opposite signs for every sufficiently small \(c\). The quotient \(-I/J\) is real analytic on a neighborhood of each closed bracket and nonconstant there by those signs. Its critical values on the finite union of those brackets form a finite set. Choose a common small rational regular value \(c\) avoiding cancellation of that coefficient. Every bracket now contains a simple zero of \(I+cJ\). One sufficiently small positive rational \(\epsilon\) gives simultaneous hyperbolic cycles by first-order return expansion. This coefficient choice gives the hyperbolic refinement at the same \(G(n,r,\mathbf m)\).

\paragraph{The nine-term quartic.}
Gasull and Santana \cite[\S1]{GasullSantana2024} attribute a nine-term quartic with 24 limit cycles to Br\'ehard, Brisebarre, Joldes and Tucker \cite[\S7.1]{BBJT2024}. For the finite calculation we use the following transcription of \cite[\S7.1, (7.2), Table~2]{BBJT2024}, with the exact inputs listed below
\[
 \begin{aligned}P&=-4y^2(y^2-11/10),\qquad Q=4xy(x^2-9/10)+\epsilon g(x,y),\\
 g&=\alpha_{00}+\alpha_{20}x^2+\alpha_{22}x^2y^2+\alpha_{40}x^4+\alpha_{04}y^4 .\end{aligned}
\]
In the displayed monomial order the five coefficients, multiplied by \(10^{20}\), are
\[
 \begin{gathered}
 -78622148667854837664,\quad 87723523612653436051,\quad 10^{20},\\
 23742713894293038223,\quad -21823846173078863753 .
 \end{gathered}
\]
The six inner energies are \(0.25,0.6084,0.7744,0.7921,0.801025,0.80766169\); the three outer energies are \(0.811801,0.8649,0.9025\). These terminating decimals specify exact rationals. The five nonzero coefficients fix \(g\), with \(\epsilon>0\) sufficiently small. The supports are \(\splitcount{2}{7}\) and \(D=4\). At the nine transcribed energies the integral below gives signs \(+,-,+,-,+,-\) on a small-oval family and \(-,+,-\) on a large-oval family. Symmetry gives four small families and two large ones, hence at least \(4\cdot5+2\cdot2=24\) isolated cycles. An earlier account explains the numerical error in a claim of 26 for a different choice of coefficients in this family \cite[\S6.2]{BrehardMahboubiPous2019}.

For the ten-term comparison marked \(^{p}\), fix one such \(\epsilon\) so that every marked endpoint has the prescribed return sign. Replace \(Q\) by \(Q+\eta x\). The \(x\) slot is vacant, the degree stays four, and the supports are \(\splitcount{2}{8}\) for \(\eta\ne0\). Return maps on the finitely many compact regular brackets depend continuously on \(\eta\). Hence one sufficiently small nonzero \(\eta\) keeps every endpoint sign and gives at least one isolated fixed point in each open bracket. This supplies the degree-preserving persistence deduction mentioned after Table~1 of Gasull and Santana \cite{GasullSantana2024}. The strict signs alone give no value for an individual return derivative.

A change of an existing coefficient gives a nearby hyperbolic quartic. In the Hamiltonian \(H=(x^2-9/10)^2+(y^2-11/10)^2\), the speed multiplier for the unperturbed field is \(y\) up to orientation. The largest mark is \(h=361/400\), so throughout every selected disk \(y^2\ge11/10-\sqrt h\ge3/20\). Varying \(\alpha_{00}\) therefore changes the Abelian integral in the direction \(\pm\oint\dd x/y\), of absolute value \(\iint y^{-2}\dd x\dd y>0\). The preceding regular-value proof gives 24 hyperbolic cycles at \(\splitcount{2}{7}\) for a nearby coefficient, and their hyperbolic continuations at \(\splitcount{2}{8}\). These multipliers belong to the perturbed coefficients. For the transcribed \(g\) the brackets give the isolated-cycle count. The attribution of these coefficients and marks is to \cite[Theorem~7.1 and Table~2]{BBJT2024}. The first five small-oval marks equal the squares of the radii in \cite[\S6.2, Table~1]{BrehardMahboubiPous2019}. That table gives enclosure widths for basis integrals. The companion archive \cite{LuReplay} gives the 160-bit interval replay for all nine marks. On small ovals take \(z=X_0+t\cos u,\ s=Y_0+t\sin u,\ t=\sqrt h\), and integrate \(t\sin(u)g(\sqrt z,\sqrt s)/(2\sqrt{zs})\) for \(0\le u\le2\pi\). On big ovals set \(z=(X_0+t)\sin^2 u,\ w=\sqrt{X_0+t}\cos u\sqrt{t-X_0+z}\), and integrate the difference of \(g(\sqrt z,\sqrt s)/\sqrt s\) at \(s=Y_0+w,Y_0-w\) against \(2\sqrt{X_0+t}\cos u\,du\) for \(0\le u\le\pi/2\). These are regular parametrizations of the two portions of the positive oval, with orientation fixed by the displayed integrands. The ball quadrature checks analyticity of each square root on complex evaluation boxes.

\subsection{Ordinary support along a two-source renewal chain}\label{sup:renewal-support}
Start with \(X=(A,B)\) of degree at most \(m\ge3\) carrying \(L\) selected hyperbolic cycles, and set \(b=m+2\). Choose \(p(t)=\prod_{i=1}^b(t-a_i)\) with distinct \(a_i\in(-1/2,1/2)\) and put \(\lambda=p(-1)/p'(-1)<0\). The two prescribed edge traces come from the polynomial field
\[
 S_b(x,y)=\bigl(p(y)+\lambda[p'(x)-p'(-1)],\quad
                 p(x)+\lambda[p'(y)-p'(-1)]\bigr).
\]
Set \(Z=S_b+\tau^{-1}(1+x)(1+y)X(x/\tau,y/\tau)\) and \(\Phi(u,v)=(2u^2-1,2v^2-1)\) \cite[Lemma~2.1, Proposition~2.2 and (21)--(26)]{LuRenewal}. For \(w=(x,y)/\tau\), the actual field of the transplant is \(\tau^{-2}\widehat Z(w)\), where \(\widehat Z(w)=\tau Z(\tau w)\to X(w)\) on compact sets in \(C^1\). The positive constant factor preserves the return map. The fold sends \(Y=(4vZ_1\circ\Phi,4uZ_2\circ\Phi)\) to \(16uv\,Z\circ\Phi\) under \(D\Phi\). Two quadrants preserve time orientation and two reverse it. Each simple root \(a\) of \(p\) satisfies the center condition \(\lambda p'(a)^2<0\) on each edge. First a perturbation \((\varepsilon v F(u),0)\) creates the vertical nests while preserving the horizontal centers. For this field a second perturbation \((0,\eta uG(v))\) creates the horizontal nests and preserves the old cycles and the vertical nests. Each of \(F,G\) has all \(b\) odd powers through degree \(2b-1\) at its admissible positive scale.

With \(U=2u^2,\ V=2v^2\), the first component is
\[
 4v\{p(V-1)+\lambda[p'(U-1)-p'(-1)]\}
 +\frac{4vUV}{\tau}A((U-1)/\tau,(V-1)/\tau)+\varepsilon vF(u).
\]
The first two supports are \((0,2j+1)\) for \(0\le j\le b\) and \((2i,1)\) for \(1\le i<b\). Each translated input position \((i,j)\) contributes \((2i+2,2j+3)\), and the odd perturbation contributes \((2i+1,1)\) for \(0\le i<b\). The second component has the exchanged edge and perturbation supports, and translated input positions \((2i+3,2j+2)\) from \(B\). The top pure powers give \(D^+=2b+1\). Let \(s_\tau(A)\) denote the number of nonzero coefficients of \(A((U-1)/\tau,(V-1)/\tau)\). The disjoint supports give
\begin{equation}\label{eq:renewal-support}
 N^+=6b+s_\tau(A)+s_\tau(B),\qquad D^+=2b+1 .
\end{equation}
Indeed the roots of \(p(t-1)\) are positive, so its coefficients and the nonconstant coefficients of its derivative are nonzero. Proposition~4.1 of the preprint gives \(4L+4b(b-1)\) selected hyperbolic cycles after these two simultaneous perturbations. From the zero field the first renewal has exactly \(6b\) terms and gives \(4b(b-1)\) cycles.

\paragraph{One doubling chain.}
Take \(X_0=0,\ m_0=3,\ m_{j+1}=2m_j+5\), so \(m_j=8\cdot2^j-5\). Write \(N_{j+1}\) for the actual term count at the next step. For \(j\ge1\), sufficiently short positive transplants at successive steps give
\[
 6(m_j+2)+2(2^j-1)(7\cdot2^j-2)
 \ \le\ N_{j+1}\ \le\ (m_j+2)(m_j+7).
\]
For the upper bound, each translated component occupies at most \(\binom{m_j+2}{2}\) positions. For the lower bound, the first top form at stage 1 is a positive multiple of \(v^{11}\). At the next step the first top form is a positive pure \(v\)-power plus a positive multiple of \(u^2v^3\) times the preceding first top form evaluated at \(u^2,v^2\). The second component has the exchanged weights. The edge derivative and both odd perturbations have lower degree, so all nonzero top coefficients stay positive. Starting at \((0,11)\), the map \((a,c)\mapsto(2a+2,2c+3)\) gives a first-component top exponent \((a_j,c_j)=(2^j-2,7\cdot2^j-3)\), with its transpose in the second component. At each position \((i,k)\) below this pair, the leading coefficient in \(1/\tau\) of the translated polynomial is a nonempty binomial sum of sign \((-1)^{m_j-i-k}\). Choosing one small \(\tau\) preserves both rectangles and the selected finite family, giving the lower bound for this seed, these binary degree choices and these successive small scales.

\paragraph{Rational coefficients and coprimality.}
Starting at the zero field, choose rational edge roots and then rational parameters in the successive admissible intervals. For the last amplitude write \(Q=Q_0+\eta uG(v)\) with \(P\) fixed. The edge trace gives \(P(0,v)=4vp(2v^2-1)\not\equiv0\). For \(c\ne0\) the odd part of \(P(u,c)/c\) is \(\varepsilon F(u)\not\equiv0\). Hence any irreducible factor of \(P\) dividing \(uG(v)\) over \(\mathbb C\) can only be \(v\), and \(Q_0(u,0)=4up(2u^2-1)\not\equiv0\). A factor of \(P\) dividing \(Q\) at two amplitudes would divide both \(uG(v)\) and \(Q_0\). Avoiding finitely many amplitudes inside the common interval gives coprime components at the stated support. The same preprint prescribes a finite Taylor jet at a point where the congruence divisor is nonzero and studies common factors forced by the congruence \cite[{\S}7.1, p.~17, and Proposition~7.3, p.~18]{LuRenewal}. In the main article, Lemma~\othernumber{lem:whole-disk-division}{3.2} prescribes an integral variation across a disk whose projected interval can contain zeros of the divisor.

\subsection{Term counts in mechanical renewal}\label{sup:mechanical-budget}
Proposition~3.3 of the mechanical renewal preprint \cite[pp.~6--10]{LuMechanical} constructs \(H(x,y)=U(x)-V(y)\) with even potentials of degree \(2n\). For \(n\ge2\) and \(r=\lfloor\log_2 n\rfloor\), the integer count \(L(n)\) in the accompanying recurrence satisfies
\[
 L(n)\ \ge\ \tfrac12 n^2r^2+\tfrac16 n^2r-\tfrac{22}{3}n^2.
\]
The proposition constructs a single polynomial form \(\omega=A\,dx+B\,dy\). On disjoint regular annuli its \(L(n)\) selected simple period zeros yield hyperbolic cycles of one field \((H_y-\mu B,-H_x+\mu A)\) of degree at most \(2n-1\), for one sufficiently small \(\mu\) \cite[Proposition~3.3 and (3.7)--(3.10)]{LuMechanical}. Its formulas also control the ordinary support. The two-coordinate fold is \(\Phi(x,y)=(x^2-a,y^2-b)\). The \(dy\) coefficient of a pulled-back form is \(2yB\circ\Phi\), and the \(dx\) coefficient is \(2xA\circ\Phi\). The added vertical \(dy\) coefficient is a sum of \(2y(c_a x)^{2p+1}R_p(y^2-b)\). The horizontal \(dx\) coefficient has the exchanged variables. Together with \(H_y=-V'(y)\) and \(H_x=U'(x)\), these formulas give odd \(y\)-exponents in the first component and odd \(x\)-exponents in the second. The exponent in the other coordinate is subject to the total degree bound.

For degree at most \(2n-1\), the exponent \(2j+1\) in one coordinate allows \(2(n-j)-1\) choices for the exponent in the other. Summing for \(0\le j<n\) gives \(n^2\) positions per component. Given \(N\), put \(n=\lfloor\sqrt{N/2}\rfloor\). For \(n\ge2\),
\[
 2n^2\le N\le 2(n+1)^2-1\le 2n(2n+1).
\]
The rightmost expression counts all component positions of total degree at most \(2n-1\). For the field supplied by the proposition the selected cycles persist in an open coefficient neighborhood. First populate the two parity supports with small coefficient changes there. Rational choices are dense. The first component has a pure \(y\) term, so \(x\) does not divide it. Varying the coefficient of \(x\) in the second component excludes at most one value for each irreducible factor of the first. This gives a coprime rational pair with \(2n^2\) terms. For any fixed \(n\ge2\) one can now fill to every \(m\) with \(2n^2\le m\le(\Delta+1)(\Delta+2)\) whenever \(2n-1\le\Delta\), by the proof of Theorem~\othernumber{thm:joint-main}{5.1}. In the floor window above, the exact recurrence gives \(L(2)=0\) at \(8\le N\le17\) and \(L(3)=4\) at \(18\le N\le31\). For \(n\ge3\) the final auxiliary birth alone has \(\lfloor n^2/4\rfloor-1\ge1\) zeros before copying. The displayed coarse bound first becomes positive at \(r=4\), or \(N=512\). Since \(0\le N-2n^2\le4n+1\) and \(2r+1\le\log_2 N<2r+3\), substitution into the bound for \(L(n)\) gives \(L(\lfloor\sqrt{N/2}\rfloor)\ge\tfrac1{16}N(\log_2 N)^2-O(N\log N)\). The existence conclusion uses that proposition and the filling argument above. The arithmetic replay checks the recurrence and term ceilings.

For a finite comparison with \(\mathcal E(N)\), count the births of that recurrence. At stage \(i\) its birth count \(K_i\) is \(i s_i^2\) less nonnegative integer losses, where \(s_i=\lfloor n/2^{r-i}\rfloor\). Counting the later copies gives
\[
 L(n)\le \sum_{i=1}^r4^{r-i}i s_i^2
 \le n^2\,r(r+1)/2\le n^4/4 .
\]
Indeed \(r(r+1)/2\le4^{r-1}\) by induction from \(r=1\), and \(2^r\le n\).
For \(n\ge2\) this upper estimate on the specified count is strictly below \((n^2-1)^2=\mathcal E(2n^2)\le\mathcal E(N)\) throughout its floor window. The elementary field uses its own degree budget.

\paragraph{Mixed terms in the archived family.}
For the displayed mechanical template one can count additional mixed terms at generic choices of the final birth coefficients. Write \(s=\lfloor n/2\rfloor\). The last birth on each axis uses the modes \(z^{2p+1}w^{2q}\,dw\) with \(0\le q<s\), \(0\le p\le n-2-2q\). Their number is
\[
 \sum_{q=0}^{s-1}(n-1-2q)=s(n-s)=\lfloor n^2/4\rfloor .
\]
Fix the shifts, nonzero amplitudes and inherited coefficients. In the vertical source, put \(Y=y^2\). For fixed \(p\), the new modes vary the coefficients of \(x^{2p+1}y^{4j+1}\) through the powers \((Y-b)^{2q}\), up to a common nonzero factor for that \(p\). The matrix \([Y^{2j}](Y-b)^{2q}\), with \(j,q\) in increasing order from \(0\) through \(\lfloor(n-2-p)/2\rfloor\), is upper triangular with diagonal one. The inherited auxiliary polynomial adds a fixed vector. The planar pullback and the Hamiltonian part of this component have even \(x\)-exponents. Thus each designated mixed coefficient is nonconstant affine in the fresh coefficients. The horizontal source gives the exchanged positions in the other component. Removing the finitely many coefficient hyperplanes leaves an open dense subset of the birth-coefficient space with at least \(2\lfloor n^2/4\rfloor\) occupied mixed component positions. This count concerns those generic last-birth choices. The transfer to exact \(N\) uses only the parity ceiling.

\paragraph{Sparse finite specializations at degree cap 31.}
In Proposition~3.3 of \cite{LuMechanical} take \(n=5\), so the successive half-degrees are \(1,2,5\). At the first global stage the auxiliary space for the current source has dimension one and requires zero simple zeros. Choose the zero current sources, hence \(\omega_1=0\) in (3.10) of that paper. The auxiliary forms for the future label \(t=2\) are chosen independently and retained. For each axis the level-zero birth for \(t=2\) has \(q=0,\ 0\le p\le2\), giving two selected simple zeros. They lift to four at the next auxiliary level. The fresh pairs from (3.5) there have \(q=0,\ 0\le p\le3\), or \(q=1,\ 0\le p\le1\). Their six periods give five new simple zeros. The two branches on each of two axes then give \(2\cdot2(4+5)=36\) simple zeros. The positive even padding in (3.9) preserves their regular collars. The reflection identities in (3.10) allow both axis amplitudes in one form and Lemma~2.3 of the preprint supplies one common small field parameter.

The final pullback of \(\omega_1\) vanishes. In each physical component (3.7)--(3.8) give only mixed powers with both exponents odd and total degree at most eight. There are ten possibilities, from \(2i+1,2j+1\) with \(i+j\le3\). The Hamiltonian adds at most five pure odd powers through degree nine. Thus 30 positions suffice in total. Populate these positions with rational coefficients in the open persistence neighborhood. The first component has a pure \(y\) term, so \(x\) does not divide it. Vary the pure \(x\) coefficient of the second component there to avoid its finitely many common-factor values. Ordinary filling then yields, in particular, 36 hyperbolic cycles at exact \(m=33,\ \Delta=31\).

For the empty seed the two-source count in Section~\ref{sup:renewal-support} is \(4b(b-1)\) at exactly \(6b\) terms and degree \(2b+1\), for integers \(b\ge5\). At a common cap of \(31\) the following choices all admit rational coprime filling by Remark~\othernumber{rem:exact-fill}{1.4} of the article.
\[
\begin{array}{c|ccc}
 b & \text{terms before filling} & \text{exact }m & \text{hyperbolic cycles}\\ \hline
 5 & 30 & 33 & 80\\
 7 & 42 & 47 & 168\\
 15& 90 & 90 & 840
\end{array}
\]
The respective degrees before filling are \(11,15,31\). At the same cap the stationary-turn choice \(n=16,d=15\) uses 33 terms and supplies 837 cycles throughout this range on filling.

\paragraph{Independent budgets.}
The sufficient term and degree conditions for the mechanical template are \(2n-1\le\Delta\) and \(2n^2\le m\le(\Delta+1)(\Delta+2)\). The largest \(n\) they admit is \(\min(\lfloor(\Delta+1)/2\rfloor,\lfloor\sqrt{m/2}\rfloor)\). Writing \(t=\min(\Delta,\sqrt{2m})\), one has \(n=t/2+O(1)\). The recurrence's leading expression becomes \(\tfrac18 t^2(\log_2t)^2+O(t^2\log t)\). For \(m\ge\lceil3\Delta/2\rceil\) this expression has smaller order than \(2\Delta^2\log_2\Delta\) when \(m=o(\Delta^2/\log_2\Delta)\), and larger order when \(m/(\Delta^2/\log_2\Delta)\to\infty\). If \(m\sim c\Delta^2/\log_2\Delta\) with fixed \(c>0\), its leading coefficient on this scale is \(c/4\). This leading-term comparison uses the full parity allowance for the recurrence. The minimum support of a realizing field remains undetermined.

Qin and Sun's preprint uses \(n=3^r,\ r\ge2\), and gives \(4r^2n^2/9-(n-1)^2\) cycles at degree at most \(4n-5\) \cite[Theorem~1.1 and {\S\S}3.2, 6.3]{QinSun2026}. In the displayed perturbation the factors \(T_n'(x)T_n'(y)T_a(x)T_b(y)\), for positive odd \(a,b\le n-2\), are integrated from 0 to \(x\) in the first variable. The resulting mixed powers are \(x^{2i}y^{2j-1}\) with \(1\le i,j\le n-1\). The Hamiltonian adds \(n\) pure odd positions in each component. Thus \(n^2+1\le m\) and \(4n-5\le\Delta\) suffice. The second component depends only on \(x\), and the leading \(y^{2n-1}\) coefficient in the first is a nonzero constant. Rational approximation and filling preserve coprimality and the selected hyperbolic cycles. The largest admissible power of three gives order \(u^2\log^2u\) with \(u=\min(\Delta,\sqrt m)\). Replicating a seed also changes its support, which must be counted again for an exact-term comparison. The least support capable of sustaining squared-logarithmic degree growth remains open.

\subsection{Which horizontal strips carry the doubled families?}\label{sup:lienard-product}
Proposition~\othernumber{prop:strip-product}{6.3} of the main article counts sign brackets on \(k=\lfloor(\Delta+1)/2\rfloor\) separate strips. After \(p\) vertical folds, the force \(V_p'\) has \(2^p\) odd powers and the perturbation \(F_p\) has at most \(2^pb-1\) nonconstant powers. At \(m=\lceil3\Delta/2\rceil\) terms, replicating \((V_p,F_p)\) directly gives leading coefficient \(1/4\). On the diagonal \((N,N-1)\), a terminal even fold doubles \(c_p\) to \(2c_p\) without new outer terms, yielding \(2^{p+1}\) odd force terms and at most \(2^pb-1\) even perturbation terms of degree at most \(2^{p+1}b-2\), which also raises the diagonal leading coefficient to \(1/4\). Han and Romanovski's property \(Z(n,m,k)\) concerns \(k\) odd-multiplicity cycles for sufficiently small perturbations of \(\dot y=X-\varepsilon F(y),\ \dot X=-g(y)\), with \(\deg F\le n+1,\ \deg g\le m\) and \(g(+\infty)=+\infty\) \cite{HanRomanovski}. In its first doubling branch \(n'=2n+1,\ m'=2m+1,\ k'=2k+n\). The sign-bracket realization uses \(n=s_i-1\) in this count.

The coordinate \(X=AT_k(x)\) has one regular inverse near each of the \(k\) simple zeros. There \(U=X^2/2\) and the confining vertical potential gives the outer ovals used in the proof. The intervening maxima have local form \(U=A^2/2-X^2/2\) in a different chart. A large outer bracket for \(X^2/2+V_p(y)\) therefore has no copy there by this change of variables. In the stationary construction of Proposition~\othernumber{prop:stationary-budget}{6.1} the vertical potential is odd, and reflection in \(y\) also supplies the maximum strips.

The horizontal replication in this proof works for any \(k\ge1\) whose strip degree and term count fit. The column \(\mathcal T\) in main Table~\othernumber{tab:fold}{2} fixes \(k=\lfloor N/2\rfloor\). For example, \((p,b,k)=(6,7,449)\) at \(N=1024\) and \((8,15,3841)\) at \(N=8192\) give \(1{,}150{,}338\) and \(116{,}036{,}610\) by the same formula \(2kc_p\). In both examples the term ceiling before filling is exactly \(N\), with degrees at most \(897\) and \(7681\), respectively. Any unused positions within the allowances may again be filled with small rational coefficients.

\subsection{Replication of an isolated-cycle seed}\label{sec:seed-replication}
A generalized Li\'enard field \(X=(y-F(x),G(x))\) of degree at most \(n\) can be replicated at a calculable term cost. We first choose a common parameter giving enough hyperbolic cycles and then determine the support of the pullback.
\begin{proposition}[Replication of a finite seed]\label{prop:seed-strips}
Suppose \(X\) as above has \(K\ge1\) distinct selected isolated limit cycles.
For every integer \(k\ge1\) and every even integer \(t\ge2\), there exists a rational polynomial field with coprime components and pure univariate terms having at least \(tkK\) simultaneous hyperbolic limit cycles. Its term count and degree are at most
\[
 k+t(n+1),\qquad \max\{2k-1,t(n+1)-1\},
\]
respectively.
\end{proposition}

\begin{proof}
Consider \(X_a=(y-F(x)-aG(x),G(x))\). On a selected orbit of period \(T\) at \(a=0\), take an analytic transversal with coordinate \(s=0\) at the orbit and let \(d(s,a)\) be the return displacement. Its parameter derivative at \((0,0)\) is nonzero. Indeed, with the starting point fixed, the parameter variation \(w\) of the flow satisfies
\[
 \frac{\mathrm d}{\mathrm dt}\det(X,w)
   =(\operatorname{div}X)\det(X,w)+\det(X,\partial_aX_a|_0),
 \qquad \det(X,\partial_aX_a|_0)=G(x)^2 .
\]
The initial determinant is zero. The last term is positive somewhere on the orbit, since \(G(x)=\dot y\) cannot vanish identically on a nonconstant periodic orbit. Integration with the positive integrating factor gives a nonzero final determinant. The derivative of the hitting time adds a vector parallel to \(X\), and conversion to the transverse coordinate divides by a nonzero determinant.

The analytic zero set of \(d\) is therefore a graph \(a=A(s)\). Its order at zero is the finite multiplicity of \(d(s,0)\), since \(d(s,a)=(a-A(s))u(s,a)\) with \(u(0,0)\ne0\). After shrinking the section, \(A'(s)\ne0\) off zero. An odd-order graph meets each sufficiently small signed level once. An even-order graph meets levels of one sign twice and levels of the other sign zero times. At each intersection \(d_s=-uA'\ne0\), so the corresponding cycle is hyperbolic. Choose disjoint collars for the \(K\) cycles and one common bound on \(|a|\). If \(o\) cycles have odd order and \(e_+,e_-\) even-order cycles open to the positive and negative sides, the larger of the two local counts is
\[
 o+2\max(e_+,e_-)=K+|e_+-e_-|\ge K .
\]
Fix the favorable sign and choose \(K\) simple cycles of that one field. Rational approximation of \(F+aG\) and \(G\) within degree \(n\) preserves all of them. Write \(F,G\) for these rational polynomials.

Choose \(k\) distinct rational nodes \(b_j\) symmetric about zero. Hermite interpolation gives a unique polynomial \(\eta\) of degree at most \(2k-1\) with \(\eta(b_j)=0,\ \eta'(b_j)=1\) at every node. The data and uniqueness give rational coefficients and \(\eta(v)=-\eta(-v)\). Thus \(\eta\) has at most \(k\) monomials. Put \(H_R(v)=R\eta(v/R)\). In \(v=Rb_j+s\) one has \(H_R(Rb_j+s)=s+O(R^{-1})\) in \(C^1\) on each fixed compact \(s\)-interval. Consequently \((H_R(v)-F(x),G(x))\) converges in each translated chart to the hyperbolic seed. Take one sufficiently large integer \(R\) for the \(kK\) compact collars, with disjoint strips for different nodes.

Next take an integer \(C>0\) with the \(x\)-projections of these collars in \((-C,C)\). Put \(\phi(u)=CT_t(u)\), where \(T_t(\cos\theta)=\cos(t\theta)\). Each of the \(t\) monotone branches in \((-1,1)\) has a smooth inverse over the selected projections. The field
\begin{equation}\label{eq:seed-fold-field}
 Y=\bigl(H_R(v)-F(\phi(u)),\,\phi'(u)G(\phi(u))\bigr)
\end{equation}
satisfies \(D(\phi,\mathrm{id})Y=\phi'(u)(H_R(v)-F(\phi(u)),G(\phi(u)))\). On each branch its forward return is conjugate to the old return or to its inverse according to the sign of \(\phi'\). A positive multiplier different from one and its reciprocal are both hyperbolic. The \(tkK\) lifted cycles lie in distinct collars of this single field.

Since \(t\) is even, \(F(\phi)\) uses at most \(tn/2+1\) even powers of \(u\), including the constant, and \(\phi'G(\phi)\) uses at most \(t(n+1)/2\) odd powers. Adding \(k\) gives at most \(k+tn+t/2+1\le k+t(n+1)\), including \(t=2\). The two component degrees give the stated maximum. The second component is a nonzero polynomial in \(u\), and the highest \(v\)-coefficient in the first is a nonzero constant, so the pair is coprime.

For a cap \(\Delta\) and \(k+t(n+1)\le m\le(\Delta+1)(\Delta+2)\) with both degrees within the cap, Remark~\othernumber{rem:exact-fill}{1.4} of the main article gives exact \(m\) and rational coprime components. For every \(m\le3\Delta+2\) one may fill within pure-power and constant positions. These include \(v,\ldots,v^\Delta\) in the first component and \(1,u,\ldots,u^\Delta\) in both, giving \(3\Delta+2\) positions that include the original support. Adding small rational coefficients there preserves the displayed coprimality argument.

\end{proof}

For a published finite seed, Huzak and De Maesschalck give \(2\lfloor(n-2)/2\rfloor+\lfloor m/2\rfloor\) hyperbolic cycles for a field \((y-F(x),G(x))\) with \(\deg F=n\ge6,\ \deg G=m\ge2\) \cite[Theorem~1.1, p.~2]{HuzakDeMaesschalck2014}. Taking both degrees equal to \(12\) gives \(K=16\). With \(k=22,\ t=2\), Proposition~\ref{prop:seed-strips} gives \(2\cdot22\cdot16=704\) cycles using at most \(22+2(12+1)=48\) terms and degree at most \(\max(43,25)=43\). Filling vacant pure powers through degree \(43\) gives exactly \(48\) terms while preserving the cycles and coprimality. Thus this count applies under the diagonal cap \(47\).

\paragraph{A conditional comparison with linear restoring force.}
Assume the seed theorem of the unrefereed preprint \cite[Theorem~1.1]{LuClassical}. For each \(n\ge14\) it asserts that, with \(G(x)=-x\), an open set of degree-\(n\) primitives \(F\) gives at least \(K=\lceil(7n-48)/2\rceil\) simultaneous hyperbolic cycles. Choose a rational seed in that open set and an integer \(C\) whose open interval \((-C,C)\) contains the horizontal projections of the selected collars. A vertical translation replaces \(F\) by \(F-F(-C)\). Choosing \(R\) large in the Hermite construction keeps the translated copies inside the same horizontal range. For \(\phi(u)=C(2u^2-1)\), the first component of \eqref{eq:seed-fold-field} then occupies at most \(k\) odd \(v\)-powers and \(n\) positive even \(u\)-powers. The second is \(-\phi'\phi\), with just \(u,u^3\). Small rational changes fill these \(k+n+2\) positions, preserve the \(2kK\) cycles, and give degree \(\max(2k-1,2n,3)\). The leading \(v\)-coefficient is a nonzero constant, ensuring coprimality. With this seed the count under the diagonal cap is
\[
 n=\lfloor(N-1)/2\rfloor,\qquad k=N-n-2,\qquad
 \mathcal C(N)=2k\left\lceil(7n-48)/2\right\rceil\quad(n\ge14).
\]
On the grid \(N=2^j,\ 6\le j\le20\), the folded count \(\mathcal F(N)\) of the article, equation~(\paperlink{main}{equation.4.7}{eq:finite-main}{4.7}), first exceeds this quadratic count at \(2^{16}\). There \(\mathcal F(N)=7{,}547{,}387{,}904\) and \(\mathcal C(N)=7{,}514{,}193{,}974\).

\section{Degree-capped specializations}\label{sup:cap}
At four terms the elementary polynomial \(\dot x=2y,\ \dot y=-2x+\mu(-y+4y^3/3)\) has degree three. Its normalized first moment is proportional to \(E(E-1)\), so one small nonzero rational \(\mu\) gives a simple return. For one, two or three terms the zero bound is due to Buzzi, Carvalho and Gasull \cite[Theorem~1.2]{BuzziCarvalhoGasull2021}.

\subsection{Two simple roots at one focus and a two-sheeted lift}\label{sup:focus}
Use the \(k=2\) truncation of Gasull and Santana's (21) \cite[pp.~11--12]{GasullSantana2024}, as written in Section~\ref{sup:sources}, with \(b=a_2\). On \(y>0\), let \(u=\log y,\ w=x-1/y\). Then
\[
 \dot u=w,\quad \dot w=g(u)+f(u)w,\quad
 g=a_1(e^u-1)+b(e^{3u}-1),\quad f=a_1e^{2u}+e^{-u}.
\]
The point \((u,w)=(0,0)\) has \(A=-g'(0)=-a_1-3b>0\) near \((a_1,b)=(-1,0)\). Define \(V(u)=-\int_0^ug\), \(z=u\sqrt{2V(u)/(Au^2)}\). This real analytic chart has \(z'(0)=1\). Dividing the field by \(z'\) gives \(z'=w,\ w'=-Az+\widetilde f(z)w\), where primes on the field now refer to the new time, and \(\widetilde f(z)=f(u(z))\,du/dz\). Write \(c_{2j}=[z^{2j}]\widetilde f\). In particular \(c_0=1+a_1\). On \(a_1=-1\), differentiation of the primitive of \(f\) and Lagrange inversion give
\[
 c_{2j}=[u^{2j}](-e^{2u}+e^{-u})
 \left(\frac{2V(u)}{Au^2}\right)^{-(2j+1)/2},\quad
 V=e^u-1-u-b((e^{3u}-1)/3-u).
\]
\hypertarget{focus-jet-coefficients}{Expansion yields} \(c_2=9b/(3b-1)\), \(c_4|_{b=0}=1/3\). Thus \((c_0,c_2)\) are local coordinates in the two parameters. 

Let \(K_j=\int_0^{2\pi}\cos^{2j}\theta\sin^2\theta\,d\theta>0\) for \(j=0,1,2\), and \(C=K_2/3\). Choose small \(\delta>0\), with \(K_0c_0=4C\delta^4,\ K_1c_2=-5C\delta^2\). In \(z=\delta Z,\ w=\sqrt A\,\delta W\) on a compact annulus about radii 1 and 2, the time-normalized system has linear part \(Z'=W,\ W'=-Z\), with damping \(\widetilde f(\delta Z)/\sqrt A\). The odd part alone gives a center: the reflection \((Z,W,t)\mapsto(-Z,W,-t)\) fixes the vertical axis, and for small \(\delta\) every trajectory in the annulus crosses it twice with nonzero speed. Its one-turn displacement is identically zero. The even part divided by \(\delta^4\) tends in \(C^2\) to \(4C/K_0-5CZ^2/K_1+Z^4/3\). To see the leading return, multiply this bounded normalized part by an independent parameter \(\eta\), integrate the first energy variation on a section, and then set \(\eta=\delta^4\). Smooth dependence of solutions and their section derivatives on a compact annulus makes the \(O(\eta^2)\) remainder uniform in \(C^1\). As \(\delta\to0\) the odd centers converge to the circle; the change of the half squared radius, divided by \(\delta^4\), therefore tends, up to the chosen orientation, to
\begin{equation}\label{eq:capped-focus-return}
 C R^2(R^2-1)(R^2-4).
\end{equation}
Its two positive roots are simple. Equivalently one may take the explicit rational family \(a_1=-1+\delta^4/6,\ b=5\delta^2/54\) for all sufficiently small positive rational \(\delta\). Indeed \(K_0=\pi, K_1=\pi/4, K_2=\pi/8\), so the same two scaled coefficients have the required limits. Both cycles persist for a rational choice of \((a_1,b)\) near the selected parameters with \(a_1b(a_1+b)\ne0\). They lie in \(x,y>0\) near \((1,1)\). Set \(x=v^2\) alone and multiply the pullback by \(2v\). Its ordinary polynomial components are
\[
 P_*=a_1v^2 y^2+b y^3-(a_1+b),\qquad Q_*=2v^3y-2v.
\]
There are three and two distinct nonzero terms, of actual degree four. On each half-plane the pullback is a diffeomorphism and the speed factor never vanishes. On the time-reversed copy the section multiplier is the reciprocal of the original one. The two disjoint copies each contain two hyperbolic cycles. This one-coordinate specialization explains the four in the capped five-term row.

For the four-term truncation take \(b=0\) and \(a_1=-1-\delta^4/24\). Then \(K_0c_0=-C\delta^4\) and \(c_2=O(\delta^4)\), since \(c_2(-1,0)=0\). The same scaled return has leading polynomial \(CR^2(R^4-1)\), with one simple positive root. Squaring both coordinates as in Section~\ref{sup:sources} lifts this cycle to four hyperbolic cycles, one per quadrant, at degree seven with support \(\splitcount{2}{2}\). Applied to the five-term truncation with two simple returns from \eqref{eq:capped-focus-return}, the same map gives eight hyperbolic cycles, two per quadrant, at degree seven with support \(\splitcount{3}{2}\).

\paragraph{Three simple cycles in the six-term focus family.}
For \(k=3\) keep the additional \(c(e^{4u}-1)\) in \(g\), with \(b=a_2,\ c=a_3\); the damping \(f\) is unchanged. On \(a_1=-1\), \(A=1-3b-4c\). In the coefficient formula above replace \(V\) by \(V-c((e^{4u}-1)/4-u)\). Expansion gives
\[
 c_2=\frac{9(b+2c)}{3b+4c-1},\qquad
 c_4\big|_{b=-2c}=-\frac{19c-1}{3(2c+1)} .
\]
At \((a_1,b,c)=(-1,-2/19,1/19)\), \(A=21/19\), \(c_6=-4/21\), and
\(\det\partial(c_2,c_4)/\partial(b,c)=6859/147\). Thus \(c_0,c_2,c_4\) are local parameter coordinates. For \(j=0,1,2\) set
\[
 K_j c_{2j}=-(4/21)K_3\,\delta^{6-2j}
 [t^j]\big((t-1)(t-4)(t-9)\big).
\]
On compact annuli in the scaled radius \(R\), the odd damping again gives a center. After division by \(\delta^6\) the even damping converges in \(C^1\) to its prescribed polynomial. The first variation about the odd center tends to its circular average, as above. The normalized energy displacement has limit proportional to \(R^2(R^2-1)(R^2-4)(R^2-9)\), with nonzero proportionality constant. All three positive roots are simple. Rational approximation inside the resulting open parameter set preserves them and all four first-component coefficients. For \(k=1,2,3\), a single square \(x=v^2\) gives \((P(v^2,y),2vQ(v^2,y))\); the two squares give the formula in Section~\ref{sup:sources}. Each sheet carries the same simple fixed points, with inverse return map on a reversed sheet. The components in all three presentations are coprime: on the torus, \(xy=1\) restricts \(P\) to \(a_1y+\sum_{j=2}^k a_j y^{j+1}-a\), a nonzero polynomial. After the square substitutions the same identity holds on each lifted component. The constant term \(-a\ne0\) and the distinct coordinate factors in the lifted components exclude common axes.

\subsection{A quartic vertical potential at six, seven and eight terms}\label{sup:quartic}
Let \(V=(y^2-1)^2,\ U=A^2T_n(x)^2,\ (\dot x,\dot y)=(V',-U'+\mu F)\), and take \(F\) odd. On each of the \(n\) horizontal inverse branches of \(X=AT_n(x)\) near its simple zeros, after multiplication by \(A\) the absolute horizontal Jacobian tends in \(C^2\) to a positive constant on compact \(X\)-intervals. This also holds for \(n=1\), when the coordinate is linear. For
\[
 L_j(y)=\sum_{k=0}^j\binom{2j+1}{2k}y^{2j+1-2k}(2-y^2)^k
\]
and \(w=y^2-1\), the difference of the drive primitive \(L_j\) at \(a=\sqrt{1+w},b=\sqrt{1-w}\) is \((a-b)^{2j+1}\). Indeed \(L_j(a)=((a+b)^{2j+1}+(a-b)^{2j+1})/2\). The negative well has the same density integral by parity. Thus, up to a common positive normalization, the first moment of \(L_j\) in the disk \(X^2+w^2<E\) begins with \(k_j E^{j+1}\), where
\[
 k_j=(2j+1)\binom{2j}{j}/(4^j(j+1)).
\]
The identity follows by differentiating the odd trace and integrating the resulting \(w^{2j}\) on a disk. For fixed large \(A\), the even part of its analytic horizontal weight equals its value at zero plus \(O(X^2/A^2)\). The leading power and normalized coefficient are unchanged. Set
\[
 F_\delta=\sum_{j=0}^{d-1}b_j L_j,\qquad
 b_j=\frac{k_{d-1}}{k_j}\delta^{d-1-j}[s^j]\prod_{i=1}^{d-1}(s-i).
\]
Now \(b_{d-1}=1\), and on \(E=\delta s\) the sum of moments divided by \(\delta^d\) converges in \(C^1\) to \(k_{d-1}s\prod_{i=1}^{d-1}(s-i)\). The analytic expansion with the scaled unit disk justifies differentiation. This yields \(2n(d-1)\) simple inner roots.

For \(E>1\) an exterior oval surrounds both vertical wells. At constant horizontal weight and \(F=L_{d-1}\) its moment has the sign, up to one fixed orientation, of
\[
 J_d(E)=\int_0^{\pi/2}L_{d-1}'(\sqrt v)\cos^2\theta\sqrt{\sqrt E-1+v}\dd\theta,\qquad
 v=(1+\sqrt E)\sin^2\theta .
\]
This follows from \(y=\sqrt{1+\sqrt E}\sin\theta\) in the area integral. \hypertarget{quartic-moment-signs}{Table~\ref{tab:capped-exterior-signs} gives rational enclosures} for \(1000 J_d(E)\). The program \path{verify_low.py} reproduces these endpoints in \path{outputs/low-N.json}.
\begin{table}[htbp]\centering\small
\caption{Rational enclosures for \(1000J_d(E)\) at the exterior test energies.}\label{tab:capped-exterior-signs}
\begin{tabular}{rrrr}\toprule
\(d\)&\(E\)&lower&upper\\\midrule\inputrows{tex/quartic-rows.tex}\bottomrule
\end{tabular}
\end{table}
Divide \([0,\pi/2]\) into 8192 equal panels, enclose sine and cosine on the whole panel in 120-bit Arb arithmetic, evaluate the exact polynomial \(L'_{d-1}(\sqrt v)\) by Horner's rule in \(v\), and sum interval ranges times panel widths. This enclosure requires no quadrature remainder assumption. First take one large integer \(A\) retaining all strict exterior signs on all strips. Then take one small rational \(\delta>0\) retaining the signs and the inner simple brackets. On every exterior annulus adding \(\gamma y\) to \(F_\delta\) changes the first integral by \(\gamma\) times a strictly positive disk area. A common small rational regular value of the quotient, as in Section~\ref{sup:sources}, gives simple exterior roots while the inner roots persist. If an ordinary odd coefficient vanishes, add a sufficiently small nonzero rational in that vacant odd slot. The leading coefficient of \(L_{d-1}\) is \(\Re(1+i)^{2d-1}\ne0\). This fills exactly the \(d\) odd powers of \(y\), preserving all simple roots through the finite perturbations. One final small rational \(\mu>0\) realizes the cycles by the main return lemma. Their collars are disjoint within each strip and across strips. The support is \(\splitcount{2}{n+d}\), the actual degree is \(\max(3,2n-1,2d-1)\). For \((n,d)=(1,2),(2,2),(2,3),(2,4)\) this proves \(3,6,10,16\) cycles at degrees \(3,3,5,7\).


For the remaining stationary entries of Table~\othernumber{tab:small}{3} in the main article, Proposition~\othernumber{prop:stationary-budget}{6.1} there uses \((n,d)=(3,4),(4,4)\), giving \(25,35\) hyperbolic cycles at degree eight.

\section{Two-sided returns across an odd stationary turn}\label{sec:stationary}
In Proposition~\paperlink{main}{theorem.6.1}{prop:stationary-budget}{6.1} of the article, subtracting a multiple of \(V'\) from the density reduces the degree of the drive while preserving its oval integrals. For the fixed field in Section~\ref{app:783-proof} we use the density before this substitution. We give its coefficients and the moment expansion used to certify its returns.

\begin{theorem}[Two-sided returns across an odd stationary turn]\label{thm:stationary-turn}
For every pair of integers $n\ge 2$ and $d\ge 4$, there exists a polynomial vector field $X=(P,Q)\in\Q[x,y]^2$ with coprime components, exact component supports $\splitcount{2}{n+d}$ and actual degree $D(X)=\max(2n-1,2d+1)$ having at least
\begin{equation}\label{eq:stationary-count}
  (2n-1)(2d-3)
\end{equation}
pairwise distinct, simultaneously coexisting hyperbolic limit cycles.
\end{theorem}

Fix integers $n\ge 2$ and $d\ge 4$, set $r:=2d-1$ (odd) and $a:=\sqrt{r/(r+2)}$, and consider
\begin{gather}
  P(x,y) = V'(y), \qquad Q(x,y) = -U'(x) + A\lambda F(y), \notag\\
  V(y) = y^r(y^2-1), \qquad U(x) = A^2 T_n(x)^2,\label{eq:st-field}
\end{gather}
where $F(y)=\int_0^y G(t)\dd t$ for an even polynomial $G\in\Q[y]$. Note that $P(x,y)=(r+2)y^{r+1}-ry^{r-1}$ has exactly two ordinary monomials, while $U'(x)=A^2(T_{2n}(x)/2)'$ occupies all $n$ odd powers $x,x^3,\dots,x^{2n-1}$. Near each of the $2n-1$ interior extrema $x_i=\cos(i\pi/(2n))$ ($1\le i\le 2n-1$, with $\delta_i=+1$ at minima and $\delta_i=-1$ at maxima), put $y=\delta_i Y$ and
\begin{gather}
  x = x_i(X) := \cos\left(\frac{i\pi}{2n} - \frac{\arcsin(X/A)}{n}\right), \notag\\
  g_i(X) := A x_i'(X) \xrightarrow[A\to\infty]{C^\infty_{\mathrm{loc}}} c_i := \frac{\sin(i\pi/(2n))}{n} > 0.\label{eq:st-chart}
\end{gather}
Because $V$ and $F$ are odd, multiplying time derivatives by $x_i'(X)>0$ transforms~\eqref{eq:st-field} on both minimum and maximum strips into $X'=V'(Y)$, $Y'=-2X+\lambda g_i(X)F(Y)$ with unperturbed Hamiltonian $h=X^2+V(Y)$. Since $V(a)<0<V(-a)$ and $V(y)=-y^r(1-y^2)$ has an odd-order stationary inflection at $y=0$ ($h=0$), for small $h\ne 0$ the bounded oval around $(0,a)$ has turning points $\ell(h)\in(-a,a)$ ($\ell(h)\to 0$ as $h\to 0$) and $u(h)>a$ ($u(h)\to 1$), and its constant-weight first energy displacement per clockwise turn is $2c_i I[G](h)$, where
\begin{equation}\label{eq:st-integral}
  I[G](h) := \int_{\ell(h)}^{u(h)} G(y)\sqrt{h-V(y)}\dd y.
\end{equation}

\begin{lemma}[Two-sided Beta expansion across an odd stationary turn]\label{lem:st-expansion}
For $j\ge 0$, set $I_j:=I[y^{2j}]$, $w:=|h|^{1/(2r)}$, $\alpha_j:=(2j+1)/r$, $\nu_j:=\alpha_j+1/2$, and $e_j:=2r\nu_j=r+4j+2$. For small $w>0$, $I_j(\pm w^{2r})$ is the sum of a common real-analytic series in signed energy $h=\pm w^{2r}$ (whose constant and linear coefficients are $S_{0j}=\frac{1}{2}\Bta(j+\frac{1}{2}+\frac{r}{4},\frac{3}{2})$ and $S_{1j}=\frac{1}{4}\Bta(j+\frac{1}{2}-\frac{r}{4},\frac{1}{2})$) and a convergent fractional series in $w^{e_j+4k}$ ($k\ge 0$) whose coefficient of $w^{e_j}$ is
\begin{equation}\label{eq:st-Cpm}
  C_j^+ = b_j\bigl(1+\sec(\pi\alpha_j)\bigr) \; (h>0), \quad
  C_j^- = -b_j\tan(\pi\alpha_j) \; (h<0), \quad
  b_j := \frac{1}{r}\,\Bta\bigl(\alpha_j,\tfrac{3}{2}\bigr).
\end{equation}
In particular, for every $j$ with $0<\alpha_j<1$ (i.e.\ $0\le j\le d-2$), $\alpha_j\ne 1/2$ because $r$ is odd, and $C_j^+ C_j^- < 0$. At $j=d-1$ both coefficients vanish, and the first nonzero fractional power on each side is $w^{e_{d-1}+4}$.
\end{lemma}

\begin{proof}
Apply the upper-end subtraction in the proof of Lemma~\paperlink{main}{theorem.6.2}{lem:stationary-moments}{6.2} of the main article to any fixed \(j\ge0\). In the coordinate \(z=y(1-y^2)^{1/r}\), each term \(z^{2i}\) in \(y^{2j}\,dy/dz\), \(i\ge j\), contributes the common upper-end primitive
\begin{equation}\label{eq:st-upper-series}
 \sum_{m\ge0}\binom{1/2}{m}\frac{h^m z_0^{2i+1+r(1/2-m)}}{2i+1+r(1/2-m)} .
\end{equation}
The half-integer denominators have absolute value at least \(1/2\) for every \(i\). The beta evaluations and normal convergence argument of that proof apply to every tail \(i\ge j\), giving \eqref{eq:st-Cpm}. For \(0<\alpha_j<1\), \(C_j^-/C_j^+=-\tan(\pi\alpha_j/2)<0\).

For \(j=d-1\) we have \(\alpha_{d-1}=1\). At both turning points \(h-V=0\), so integration of \(\frac{\mathrm d}{\mathrm dy}(h-V)^{3/2}\) gives
\[
 0=\int_{\ell(h)}^{u(h)} V'(y)\sqrt{h-V(y)}\dd y=(r+2)I_d(h)-rI_{d-1}(h).
\]
Since \(1<\alpha_d=1+2/r<3/2\), both \(C_d^+\) and \(C_d^-\) in \eqref{eq:st-Cpm} are nonzero. The first fractional power of \(I_{d-1}\) is therefore the one of \(I_d\), with \(e_d=e_{d-1}+4\).
\end{proof}

\begin{proof}[Proof of Theorem~\ref{thm:stationary-turn}]
We use the density of Proposition~\othernumber{prop:stationary-budget}{6.1} of the main article \emph{before} replacing its highest power. To specify also the fixed field of Section~\ref{app:783-proof}, put \(b=d-2,\ R_{ij}=S_{ij}/S_{ib}\), and let \(L(s,t)\) be the unique combination of \(y^{2b},y^{2d}\) whose normalized constant and linear coefficients are \(s,t\). The control columns are independent by Lemma~\othernumber{lem:stationary-moments}{6.2}. The densities \(L(1,0),L(0,1)\) and \(y^{2j}+L(-R_{0j},-R_{1j})\) for \(0\le j<b\) have respective leading orders \(0,2r,e_j\). Their coefficient pairs are \((S_{0b},S_{0b})\), \((S_{1b},-S_{1b})\), and \((C_j^+,C_j^-)\). Order the densities \(D_i\) by increasing exponents \(\tilde e_i\), denote their coefficients by \(c_{i,\pm}\), and set
\begin{equation}\label{eq:st-choice}
 G(y)=\sum_{i=0}^{d-1}\sigma_i q^{(B-\tilde e_i)^2}D_i(y),
 \qquad B>\tilde e_{d-1},\quad \sigma_i=(-1)^i\sgn(c_{i,+}).
\end{equation}
The adjacent-scale argument in the proof of that proposition gives \(d-1\) positive and \(d-2\) negative simple brackets for the limiting integral of this very \(G\). Each \(j<b\) occurs in one basis column. Both control powers have a nonzero coefficient of \(q^{B^2}\) from \(L(1,0)=(y^{2d}-R_{1d}y^{2b})/(R_{0d}-R_{1d})\). Thus all \(d\) powers occur for small rational \(q>0\). Equation~\eqref{eq:st-chart} gives \(C^1\) convergence of the weighted moments to \(c_i\) times these moments on the regular brackets. One large integer \(A\) and one small rational \(\lambda>0\) work on all \(2n-1\) strips. The drive now has degree \(2d+1\). The highest \(x\)-coefficient is constant and nonzero and \(P\) depends only on \(y\), giving coprimality and the asserted resources.
\end{proof}

\section{A specified rational field with 783 hyperbolic cycles}\label{app:783-proof}
At \(N=32\), Proposition~\paperlink{main}{theorem.6.1}{prop:stationary-budget}{6.1} of the article gives 783 cycles at degree 30 for sufficiently separated parameters. Here we fix the parameters in the unreduced family of Theorem~\ref{thm:stationary-turn} and verify its returns at degree 31.

\begin{proposition}[Explicitly specified $783$-cycle rational vector field]\label{prop:783-witness}
Fix $(n,d)=(15,15)$, $r=29$, $B=96$, $q=2^{-64}$, $A=2^{10{,}000{,}000}$, and $\lambda=2^{-20{,}000{,}000}$. Let $G(y)$ be the explicit rational even polynomial~\eqref{eq:st-choice} of degree $30$, $F(y)=\int_0^y G(t)\dd t$, $U(x)=A^2 T_{15}(x)^2$, and
\[
  P(x,y) = 31y^{30} - 29y^{28}, \qquad
  Q(x,y) = -U'(x) + A\lambda F(y).
\]
Then $X=(P,Q)\in\Q[x,y]^2$ is coprime, has exact component supports $\splitcount{2}{30}$ ($N(X)=32$) and actual degree $D(X)=31$, and possesses at least $783=29\times(14+13)$ distinct simultaneous hyperbolic limit cycles. In each of the $29$ disjoint horizontal strips, the local return section is parametrized by $h=X^2+V(Y)$. At least one cycle meets each of the $14$ return intervals in $h>0$ and each of the $13$ in $h<0$.
\end{proposition}

In the coordinates of \eqref{eq:st-chart}, the physical Hamiltonian \(H=U(x)+V(y)\) equals \(h\) on a minimum strip and \(A^2-h\) on a maximum strip. For the perturbed return, \(h\) specifies the starting point on the section.

We prove Proposition~\ref{prop:783-witness} at the fixed rational parameters $n=d=15$, $r=29$, $B=96$, $q=2^{-64}$, $A=2^{10^7}$, and $\lambda=2^{-2\cdot 10^7}$, distinguishing explicitly between finite rational/interval enclosures verified computationally and uniform complex-analytic estimates holding across all ovals in the $27$ brackets.

\paragraph{1. Finite rational and outward-rounded interval enclosures.}
Here $\tilde{e}_0,\dots,\tilde{e}_{14}$ equal $0,31,35,39,43,47,51,55,58,59,63,67,71,75,79$. Exact rational polynomial arithmetic reconstructs $D_0,\dots,D_{14}\in\Q[y]$, $G(y)$, $F(y)$, and $-U'(x)$, confirming that $P$ has $2$ nonzero terms, $Q$ has $15$ odd $x$-powers and $15$ odd $y$-powers ($N(X)=32$, $D(X)=31$, $\gcd(P,Q)=1$), $\sum_k|[y^k]D_i|<2^{40}$ ($0\le i\le 14$), and $\sum_k|[y^k]F|<2^{60}$. Outward-rounded Arb ball arithmetic verifies $2^{-12}<|c_{i,\pm}|<2^{12}$ ($0\le i\le 14$) and $c_{i,+}c_{i,-}<0$ ($1\le i\le 14$). For each of the $14$ sign-exchanging adjacent pairs on $h>0$ and $13$ pairs on $h<0$ ($i\ge 1$), set $\Delta:=\tilde{e}_{i+1}-\tilde{e}_i\ge 1$, $t_0:=|c_{i,\pm}/c_{i+1,\pm}|^{1/\Delta}$, $w_0:=q^{2B-\tilde{e}_i-\tilde{e}_{i+1}}t_0$, $w=w_0 v$, and $N_{\mathrm{scale}}(v):=\sigma_i c_{i,\pm}q^{(B-\tilde{e}_i)^2}(w_0 v)^{\tilde{e}_i}$. Define explicit dyadic rational endpoints $w_-:=2^{-80}q^{2B-\tilde{e}_i-\tilde{e}_{i+1}}\lfloor 2^{80}t_0(1-\frac{1}{5\Delta})\rfloor$ and $w_+:=2^{-80}q^{2B-\tilde{e}_i-\tilde{e}_{i+1}}\lceil 2^{80}t_0(1+\frac{1}{5\Delta})\rceil$, and put $[v_-,v_+]:=[w_-/w_0,w_+/w_0]$. Outward-rounded Arb evaluation certifies that on all $27$ brackets: (i) $[v_-,v_+]\subset(1-\frac{1}{4\Delta},1+\frac{1}{4\Delta})$ and the $w$-intervals on each side are pairwise disjoint inside $(2^{-10998},2^{-2002})$, so $H_*:=2^{-10^6}<|h|<2^{-100}$ and $|N_{\mathrm{scale}}(v)|>2^{-2\cdot 10^6}$; and (ii) the normalized leading-power sum $L_{\mathrm{lead}}(v):=\sum_{k=0}^{14}\frac{\sigma_k c_{k,\pm}}{\sigma_i c_{i,\pm}}q^{(B-\tilde{e}_k)^2-(B-\tilde{e}_i)^2}(w_0 v)^{\tilde{e}_k-\tilde{e}_i}$ satisfies $L_{\mathrm{lead}}(v_-)>+\frac{1}{8}$, $L_{\mathrm{lead}}(v_+)<-\frac{1}{8}$, and $L_{\mathrm{lead}}'(v)<-\frac{\Delta}{2}$ throughout $[v_-,v_+]$.

\paragraph{2. Uniform holomorphic moment remainder on $|w|\le 1/8$ (analytic proof).}
Write $\Psi_i^\pm(w):=I[D_i](\pm w^{58})$. On $|y|=3/4$, $|y(1-y^2)^{1/29}|\ge\frac{3}{4}(\frac{7}{16})^{1/29}>\frac{1}{2}$, so Rouch\'e's theorem yields an analytic left-turn inverse $y=\psi(z)$ on $|z|\le 1/2$ with $|\psi(z)|\le 3/4$ and $\psi(1/4)\in(1/5,1/3)$. Cauchy's estimate on $|z|\le 3/8$ gives $|\psi'(z)|\le\frac{3/4}{1/2-3/8}=6\le 8$ and $|[z^{2m}]\psi^{2j}\psi'|\le 8(8/3)^{2m}$. Taking $z_0=1/4$, for $|w|\le 1/4$ ($|h|\le 2^{-116}$, $|h|/z_0^{29}\le 2^{-58}$), the series~\eqref{eq:st-upper-series} is bounded by $4z_0^{2m+31/2}$, while $|b_m|\le\frac{1}{2m+1}$ and $|\cos(\pi\frac{2m+1}{29})|\ge\sin\frac{\pi}{58}>\frac{1}{29}$ bound the fractional constants by $30$, so the left-cut integral is bounded by $\sum_{m\ge 0}8(8/3)^{2m}(4\cdot 4^{-2m-31/2}+30\cdot 4^{-31-4m})<1$. On the middle real segment $[\psi(1/4),v_0]$ with $v_0:=1-2^{-16}$, the unimodal function $-V(y)=y^{29}(1-y^2)$ exceeds $\min(-V(1/5),-V(v_0))>2^{-75}\gg 2^{-116}\ge|h|$, giving a holomorphic square root of modulus $<2$. Near the right turn $y=1$, $V(1)=0$, $V'(1)=2$, and $|V''(y)|<2^{12}$ on $|y-1|\le 2^{-12}$, so on $|y-1|=2^{-12}$ and $|v|\le 2^{-13}$ one has $|V(y)-2(y-1)|+|v|\le 2^{-13}+2^{-13}<2^{-11}=|2(y-1)|$; Rouch\'e's theorem therefore yields a holomorphic right-turn inverse $\phi_{\mathrm{right}}(v)$ on $|v|\le 2^{-13}$ with $|\phi_{\mathrm{right}}'(v)|\le 1$. Since $2^{-17}<-V(v_0)<2^{-14}$, parametrizing $[v_0,u(h)]$ by $v=V(v_0)+t(h-V(v_0))$ ($0\le t\le 1$) bounds $|I_j(\pm w^{58})|\le 2^{12}$ on $|w|\le 1/4$. Multiplying by $\sum_k|[y^k]D_i|<2^{40}$ gives $|\Psi_i^\pm(w)|\le 2^{52}<2^{100}$ on $|w|\le 1/4$, and Cauchy's estimate on $|w|\le 1/8$ ($\tilde{e}_i\le 79$) yields
\begin{equation}\label{eq:app-psi-rem}
  \bigl|\Psi_i^\pm(w)/w^{\tilde{e}_i} - c_{i,\pm}\bigr| \le 2^{100}\cdot 4^{79}\cdot 8|w| < 2^{270}|w| \qquad (|w|\le 1/8).
\end{equation}
On the complex disk $|v-1|\le\frac{1}{2\Delta}$ ($2^{-11000}<|w|<2^{-2000}$),~\eqref{eq:app-psi-rem} bounds the difference $|I[G](\pm w^{58})/N_{\mathrm{scale}}(v)-L_{\mathrm{lead}}(v)|$ by $2^{-1700}$ (and its $v$-derivative on $|v-1|\le\frac{1}{4\Delta}$ by $2^{-1690}$ via Cauchy), so $I[G]/N_{\mathrm{scale}}$ inherits strict margins $>+\frac{1}{8}-2^{-1700}$, $<-\frac{1}{8}+2^{-1700}$, and slope $<-\frac{\Delta}{2}+2^{-1690}$ on every $[v_-,v_+]$.

\paragraph{3. Complex angular ODE and full nonlinear return (analytic proof).}
In clockwise polar coordinates $(X,Y)=(\xi\cos\phi,a-\xi\sin\phi)$ around $(0,a)$, on every reference oval ($H_*<|h|<2^{-100}$) one has $2^{-16}<\xi<4$, $|X|<2$, $-1/4<Y<5/4$, and angular speed numerator $J:=2X^2+(Y-a)V'(Y)=2X^2+31Y^{28}(Y-a)^2(Y+a)>2^{-40}H_*$. For $f(\xi):=X^2+V(Y)$, $f'(\xi_0)=J/\xi_0>2^{-43}H_*$ and $|f''|<2^{100}$ on $|\delta\xi|\le 1/4$, so Rouch\'e's theorem on $|\delta\xi|=\beta:=2^{-160}H_*$ yields a holomorphic implicit radius $\xi(\phi,h')$ for $|h'-h|\le\eta:=2^{-220}H_*^2$ with $|J|>2^{-62}H_*$, $|X|,|Y|<3$, $|g_i|<2$, $|g_i-c_i|<2^6/A$, $|V'|<2^{56}$, and $|F|<2^{111}$. Along the perturbed angular ODE
\[
  \frac{\dd h}{\dd\phi} = \frac{\lambda\xi^2 V'(Y)g_i(X)F(Y)}{J - \lambda X g_i(X)F(Y)},
\]
for $|\lambda|\le R:=2^{-600}H_*^3$ we have $|\dd h/\dd\phi|\le 2^{250}|\lambda|/H_*$ and $2\pi\cdot 2^{250}R/H_*<\eta/4$, so starting at $|h'-h|\le\eta/2$ yields a holomorphic one-turn displacement $\mathcal{D}(h',\lambda)$ of modulus $<1$. Cauchy's tail bound in $\lambda$ ($2\lambda/R^2$), replacement of $g_i$ by $c_i$ ($2^{260}/(A H_*)$), and Cauchy's $h'$-derivative factor ($4/\eta$) give, at $\lambda=2^{-2\cdot 10^7}$ and $A=2^{10^7}$,
\[
  \left|\frac{\mathcal{D}(h,\lambda)}{2c_i\lambda} - I[G](h)\right| < 2^{-3{,}000{,}000}, \qquad
  \left|\partial_h\!\left(\frac{\mathcal{D}(h,\lambda)}{2c_i\lambda} - I[G](h)\right)\right| < 2^{-3{,}000{,}000}.
\]
Dividing by $|N_{\mathrm{scale}}(v)|>2^{-2\cdot 10^6}$ leaves normalized value error $<2^{-10^6}$. Since $|\dd h/\dd v|=58|h/v|<120$ and $|N'_{\mathrm{scale}}/N_{\mathrm{scale}}|=\tilde e_i/|v|<106$, the product rule bounds the normalized slope error by $(120+106)2^{-10^6}<2^{-100}$ on all $27$ brackets in each of the $29$ disjoint horizontal strips, certifying $29\times 27=783$ simultaneous hyperbolic limit cycles.

\paragraph{4. Supplementary whole-oval quadrature check and rigorous right-turn root enclosure.}
As an independent numerical cross-check of Lemma~\ref{lem:st-expansion} on both sides of $h=0$, we evaluate $I[D_i](h)$ directly by $1000$-bit outward-rounded Arb integration at $\ell=\pm 2^{-16}$ ($h=V(\ell)$). Writing $u(h)=1+z_*$ for the unique root of $V(y)=h$ in $(a,\infty)$, exact rational evaluation first verifies endpoint sign-straddling $V(1-|h|)<h<V(1+|h|)$ together with strict monotonicity $V'(1+I_0)>1$ across $I_0:=[-|h|,|h|]$. For $T(z):=h/((1+z)^{29}(2+z))$, ball arithmetic verifies strict self-map inclusion $T(I_0)\subset\operatorname{int}(I_0)$, so every iterate $I_m:=T^m(I_0)$ encloses $z_*$ by inclusion monotonicity ($z_*=T(z_*)\in I_m$); moreover, on the final narrow enclosure $[u_{\mathrm{lo}},u_{\mathrm{hi}}]$ around $1+I_5$, interval evaluation also verifies $V'([u_{\mathrm{lo}},u_{\mathrm{hi}}])>1$ on the hull and strict opposite residual signs $V(u_{\mathrm{lo}})-h<0<V(u_{\mathrm{hi}})-h$ at its endpoints before factoring $h-V(y)=(y-\ell)(u(h)-y)B(y)$.


\clearpage
\section{Bounds for the loss at one dyadic budget}\label{sup:reproduction}
In equation~(\paperlink{main}{equation.4.7}{eq:finite-main}{4.7}) of the article take \(N=2^b\) with \(b=2^{128}\). Then \(r=b-7,\ R=N/128,\ X_*=T=N/2\). For this folded count the normalized loss satisfies
\begin{equation}\label{eq:single-deficit}
 \frac{137}{4}<\delta_N:=b-\frac{\mathcal F(N)}{N^2}\le\frac{1503}{32}.
\end{equation}
For \(j=r-k\) we have \(M_k=2^{j+5}\). Put \(w_k=Y_k/2^{k+1}\) and let \(B_k^*\) be the maximum over the permitted prefixes at that stage. Dividing the stage sum by \(N^2\) gives
\[
 \delta_N=7+\sum_{j=1}^r\left(1-\frac{w_k B_k^*}{4q_j M_k}\right).
\]
Each summand is nonnegative since \(w_k\le1\) and \(B_k^*\le4q_jM_k\). Retain \(j=1,\ldots,128\) and keep only block indices \(i\le20\) in the product for \(w_k\). This can only increase the retained weights and gives a lower bound for the loss. The resulting finite rational sum exceeds \(137/4\) (\path{verify_transfer.py}, function \texttt{deficit\_certificate} \cite{LuReplay}).

For the upper bound, \(p=d=N/4\) is divisible by \(2^{k+1}\) at every stage, so \(M_k=T/2^{k+1}\). Apply the stage inequality~(\paperlink{main}{equation.5.1}{eq:Bq-linear-lower}{5.1}) and the estimates in the proof of Theorem~\paperlink{main}{theorem.1.1}{cor:linear-cap}{1.1} in Section~\othernumber{sec:sum}{5} of the article. The seven omitted levels, branch loss, two sums involving \(q_j\), and two errors proportional to \(R\) give respectively
\[
 7+\frac52+\frac{3\cdot22}{4}+\frac{41}{2}+\frac{24+36}{128}=\frac{1503}{32}.
\]

\section{A fifth-degree fold and an obstruction to prescribing a primitive}\label{sec:quintic}

The asymmetric fold \(J(z)=z^4(1+z)\) illustrates why Lemma~\othernumber{lem:whole-disk-division}{3.2} (applied below with \((A,B)=(H,B)\)) varies the potential along with the drive. Their coupled variation realizes a limiting primitive of degree thirteen with a drive correction of degree eleven. With the potential fixed, corrections of that degree cannot converge to the prescribed primitive. A second inherited drive on the same disks obstructs polynomial realization of that specified primitive at every degree. Subtraction of the two branch values on a level is the balance of Grau, Ma\~nosas and Villadelprat \cite[{\S}1, before Theorem~A]{GrauEtAl2011}. The example has $N=27$ and $D=29>26=N-1$. At the same term count Proposition~\othernumber{prop:stationary-budget}{6.1} gives $529$ hyperbolic cycles at degree $26$.

\begin{proposition}\label{prop:quintic}
Let $z=y^2-1$, $J(z)=z^4(1+z)$, and take the inherited drive $J$ and potentials
\[
 U(x)=\tau^{12}x^2,\qquad \bar V=A(J),\qquad
 A(t)=t^3/3+\tau^4 t^2/2 .
\]
For every rational $0<\tau\le2^{-40}$, polynomial changes $\lambda H(z)$ of $\bar V$ and $\lambda B(z)$ of the drive, with $\deg H\le9$, $\deg B=11$, give a simple zero of the area moment in each of the two rows near $y=\pm1$, for one common small positive rational $\lambda$. The limiting primitive in $Z=z/\tau$ has degree $13$. At $\tau=2^{-40}$, a field of the form (\othernumber{eq:separable-template}{1.1}), with rational coefficients, coprime components, exact term counts $\splitcount{15}{12}$ and degree $29$, realizes the two corresponding hyperbolic cycles. Their return multipliers are reciprocal.
\end{proposition}

\begin{proof}
In the scaled chart the potential $A(J(\tau Z))/\tau^{12}$ tends on compact sets with all derivatives to
$\Psi(Z)=Z^8/2+Z^{12}/3$. Use energies $E\in[3/4,5/4]$ for $x^2+\Psi(Z)$ and take the levels $1/2,3/2$ for the inner core and outer boundary of Lemma~\othernumber{lem:disk-core}{3.3}. These disks lie in $|Z|<2$. Both physical charts are regular there for small $\tau$. On each side of $z=0$, $J$ is monotone locally, since $J'=z^3(4+5z)$. Also $A'(t)>0$ for $t>0$. The two endpoints of each nonempty vertical slice have the same value of $J$, so the inherited drive has zero moment on each disk.

Choose $S(0)=0$ and $S'=\Psi+(3/2)Z\Psi'-1$. With $\mathcal A(E)$ the area of $\{x^2+\Psi(Z)<E\}$, integration first in $x$ gives
\begin{equation}\label{eq:quintic-moment}
 \iint_{x^2+\Psi<E} S'(Z)\,\dd x\dd Z=(E-1)\mathcal A(E).
\end{equation}
Indeed the integral of $\partial_Z[Z(E-\Psi)^{3/2}]$ between the turning points vanishes. After $Z=E^{1/8}u$, $\mathcal A(E)/E^{5/8}$ decreases with $E$ since the coefficient of $u^{12}$ in the scaled potential increases. Thus $0\le E\mathcal A'(E)\le5\mathcal A(E)/8$. The right side of \eqref{eq:quintic-moment} has opposite endpoint signs and derivative at least $19\mathcal A(E)/24$ on the bracket.

Here $D_0=J(J+\tau^4)$, $E_0=1$. The factor $C=z^4$ includes the full multiplicity of the root inside the disk, leaving $D_{\rm reg}=D_0/C=(1+z)(J+\tau^4)$. Prescribe
\[
 N_0=z^4\tau^4\bigl(1+(z/\tau)^4\bigr)S(z/\tau).
\]
After cancellation the effective primitive is
\[
 \left.\frac{N_0}{D_0}\right|_{z=\tau Z}
 =\frac{(1+Z^4)S(Z)}{(1+\tau Z)\{1+Z^4(1+\tau Z)\}} .
\]
It converges to $S$ in $C^2([-2,2])$. For $\tau\le1/4$ the denominator is at least $1/2$ on this interval. Omitting $C$ would leave a pole of order three at $z=0$. The division gives $C\mid H$ and $(\deg D_0,\deg N_0)=(10,21)$, so $\deg H\le9$ and $\deg B=11$. Indeed degree $11$ is necessary for this numerator even if $\deg H$ is allowed to reach the inherited potential degree $15$, since $D_0B-H=N_0$. A drive change alone of degree at most $11$ cannot converge under this rescaling to $S$ of degree $13$, by closedness of the space of polynomials of degree at most $11$ on a nontrivial interval.

Here one can fix the scale explicitly. Write $u=1+\tau Z$, $w=Zu^{1/4}$ and $\theta=Z^4/(1+Z^4)$. The scaled energy becomes exactly $x^2+\Psi(w)$. The effective primitive in $w$ is $T(w)=K S(Z)$, where $K^{-1}=u(1+\theta(u-1))$. For $\tau\le1/8$, $|Z|\le2$ gives $3/4\le u\le5/4$. The chart covers $|w|\le3/2$ since $|w(\pm2)|\ge 2(3/4)^{1/4}>3/2$, and $\Psi(3/2)>3/2$ puts the entire outer disk inside this chart. In the following bounds primes on \(w,K,\theta\) denote derivatives with respect to \(Z\). The mean value theorem gives $|w-Z|\le2\tau$, $|w'-1|\le2\tau$ and $w'\ge3/4$. Also $K^{-1}\ge9/16$, $|K^{-1}-1|\le9\tau/2$. Using $|\theta'|\le32$ gives $|(K^{-1})'|\le83\tau$, hence $|K-1|\le9\tau$ and $|K'|\le512\tau$. From the displayed $S'$ one has $|S|,|S'|,|S''|\le2^{19}$ on $[-2,2]$. On the outer disk these bounds give
\[
 |T'(w)-S'(w)|\le
 \frac{2^{19}(512+9+2+2)\tau}{3/4}<2^{30}\tau .
\]
Indeed differentiate $KS$ and compare $S'(Z)$ with $S'(w)$ before dividing by $w'$. At $\tau\le2^{-40}$ the error in the moment is at most $2^{-10}\mathcal A(E)$, and its energy derivative has error at most $2^{-10}\mathcal A'(E)$ by the coarea formula on the regular boundary. These bounds preserve both signs and the slope, with a common orientation reversal on the negative row.

For each such $\tau$, the first variation (\othernumber{eq:shape-derivative-green}{3.5}) and the vanishing inherited moments give both simple zeros at physical energies $\tau^{12}E$ for one sufficiently small positive rational $\lambda$. Subtract the constant drive value at $y=0$. At $\tau=2^{-40}$ the family already contains every required power of $v=y^2$ for all but finitely many $\lambda$. To check the coefficients, first evaluate them at $\lambda=1$. The products over $i=1,\ldots,15$ and $i=1,\ldots,11$ of the $v^i$ coefficients of $A(J(v-1))+H(v-1)$ and $J(v-1)+B(v-1)$ are respectively $71$ and $7$ modulo $101$. These short divisions can be reproduced over $\mathbb F_{101}$ by setting $\tau=87$, taking $H=(-N_0)\bmod D_0$ and $B=(N_0+H)/D_0$. Consequently every required coefficient is a nonzero polynomial in $\lambda$. With $\tau=2^{-40}$ fixed, choose a positive rational $\lambda$ in the common interval above, avoiding the finitely many coefficient roots. Lemma~\othernumber{lem:hyperbolic-return}{1.2} then gives an interval $0<\mu<\mu_0(\tau,\lambda)$ valid on both collars. Choose any rational \(\mu\) in this interval. At the fixed scale \(\tau=2^{-40}\), this one rational pair \((\lambda,\mu)\) gives one hyperbolic cycle near \(y=1\) and one near \(y=-1\). The resources are the $15$ odd powers through $y^{29}$ in the first component, $x$ and the $11$ positive even powers through $y^{22}$ in the second. The coefficient of $x$ is $-2\tau^{12}$, so the two components are coprime. Reflection in $y=0$ reverses time for this field. It exchanges the two collars and conjugates a return map to the inverse return on the other, giving the reciprocal multipliers.
\end{proof}

\paragraph{A failure with the same fold and the same disks.}
Replace the inherited drive $J$ by $J^2/2+\tau^4J$, keeping $N_0$ and the potential fixed. Its derivative in $J$ is now $E_0=J+\tau^4$. The inherited moment still vanishes on every disk, but $\Gamma=\gcd(D_{\rm reg},E_0)=J+\tau^4$ has degree five. For $0<\tau\le1/4$ this polynomial has a real root $\xi\in(-2,-1)$. There $R=(\tau^4+z^4)S(z/\tau)$ is nonzero, since
\[
 S(u)=u\bigl(-1+13u^8/18+19u^{12}/39\bigr)
\]
has no zero for $|u|\ge2$. Adding a constant to the primitive changes the numerator by a multiple of $D_{\rm reg}$. Since $\Gamma\nmid R$, the specified primitive has no polynomial realization at any degree, even up to a constant. Other primitives may give the same integrated moment or other cycles.

The quintic has prime degree and local order four at the origin. Its transfer uses the common value of $J$ on the two slice endpoints and the nonvanishing of the remaining factor on the disk. The branch multiplicities and recurrent critical values used in Theorem~\othernumber{thm:finite-folded}{4.1} require the separate quadratic construction of Lemma~\othernumber{lem:recurrent-folds}{3.1}.


\bibliographystyle{plain}
\bibliography{references}
\end{document}
